\section{The arc Hilbert-Poincar\'{e} series}\label{sec:section_3}

In this section we introduce and discuss the Hilbert-Poincar\'{e} series of the arc algebra of a (not necessarily reduced nor irreducible) algebraic variety $X$, focussed at a point $\mathfrak{p}$ of $X$. Since this will be a local invariant, we may restrict ourselves to $X$ being a closed subscheme of affine space. For generalities about Hilbert-Poincar\'{e} series of graded algebras we refer to the appendix (Section~\ref{appendix}).\\

As above, let $k$ be a field of characteristic zero. Although for most of the statements it is not necessary that $k$ is algebraically closed, we will assume it for convenience. Let $X$ be a subscheme of affine $n$-space over $k$, defined by some ideal $I$ in $k[x_1,\ldots , x_n]$. We define a grading on the polynomial ring $k[x_j^{(i)};1\leq j \leq n, i\in  \mathbb{N}]$ by putting the weight of $x_j^{(i)}$ equal to $i$. We prefer to use the terminology `weight' instead of `degree' here, in order not to confuse with the usual degree. It is easy to see that the ideal $I_{\infty}$ of $k[x_j^{(i)};1\leq j \leq n, i\in  \mathbb{N}]$ defining the arc space $X_{\infty}$ is homogeneous (with respect to the weight) and hence the arc algebra $J_{\infty}(X)$ is graded as well (this follows for instance from Proposition~\ref{prop:equations_by_deriving}). Similarly the jet algebras $J_m(X)$ are graded. Let $\mathfrak{p}$ be any point of $X$ and denote by $\kappa(\mathfrak{p})$ the residue field at $\mathfrak{p}$. After identifying $J_0(X)$ with the coordinate ring of $X$ we obtain a natural map from $J_0(X)$ to $\kappa(\mathfrak{p})$. 

\begin{Def}\label{def:arc_Hilbert}
We define the \emph{focussed arc algebra 
of $X$ at $\mathfrak{p}$} as
\[ J_{\infty}(X) \otimes_{J_0(X)} \kappa(\mathfrak{p}) \]
and we denote it by $J_{\infty}^{\mathfrak{p}}(X)$. Analogously, we define the \emph{focussed jet algebras of $X$ at $\mathfrak{p}$} by
\[ J_{m}^{\mathfrak{p}}(X):= J_{m}(X) \otimes_{J_0(X)} \kappa(\mathfrak{p}). \]
Using the above grading, we write $ \HP_{J_{\infty}^{\mathfrak{p}}(X)}(t)$ respectively $\HP_{J_{m}^{\mathfrak{p}}(X)}(t)$ for their Hilbert-Poincar\'{e} series as graded $\kappa(\mathfrak{p})$-algebras. We call this the \emph{arc Hilbert-Poin\-ca\-r\'{e} series at $\mathfrak{p}$} respectively the \emph{mth jet Hilbert-Poincar\'{e} series at $\mathfrak{p}$}.
\end{Def}

In fact, the focussed arc algebra at a point $\mathfrak{p}$ is the coordinate ring of the scheme theoretic fiber of the morphism $\psi_0 : X_{\infty} \to X$ over $\mathfrak{p}$. Note that the weight zero part of $J_{\infty}^{\mathfrak{p}}(X)$ or $J_{m}^{\mathfrak{p}}(X)$ is always a one-dimensional $\kappa(\mathfrak{p})$-vector space. \\

In the special case that $X$ is a hypersurface given by a polynomial $F\in k[x_1,\ldots,x_n]$ with $F(0)=0$, this boils down to the following. We define $F_0$ to be $F$ in the variables $x_1^{(0)},\ldots , x_n^{(0)}$. Then we put $F_1:= DF_0, F_2:=DF_1,\ldots $, where $D$ is the derivation from Section~\ref{sec:section_2}.
The arc algebra $J_{\infty}(X)$ is given as the quotient of $k[x_j^{(i)}; 1\leq j\leq n, i\in \mathbb{N}]$ by $(F_0,F_1,\ldots )$ (see Proposition~\ref{prop:equations_by_deriving}). And $J_{\infty}^0(X)$ is the quotient of $k[x_j^{(i)}; 1\leq j\leq n, i\geq 1]$ by $(f_0,f_1,\ldots )$, where $f_i$ is $F_i$ evaluated in $x_1^{(0)}= \cdots = x_n^{(0)} = 0$.

\begin{Rem}
Besides that the arc Hilbert-Poincar\'{e} series is very natural to look
at when working with arc spaces, its introduction is motivated by the
following. If $X$ is for instance a hypersurface then the ideal $I_m^0 = (f_0,f_1,\ldots )$ defining the fiber over the origin of the $m$th jet space of $X$ is generated by polynomials depending only on a subset of the variables
of the polynomial ring $k[x_j^{(i)}; 1\leq j\leq n, i\geq 1]$. Heuristically,
for a given $m \in \mathbb{N}$, the more $X$ is singular, the less
variables appear in $I_m^0$. So this series was meant as a
Hironaka type invariant of the singularity (see
\cite{berthomieu_hivert_mourtada}), that is a kind of measure of the
number of variables appearing in $I^{0}_m$.

Note also that since the jet spaces are far from being equidimensional
in general (see
\cite{mourtada_plane_branches,mourtada_toric_surfaces}), the jet
algebras have a big homological complexity, what makes it difficult to
compute the series introduced above.
\end{Rem}

\begin{Rem}\label{rem:graded_structure}
We can define the graded structure on $J_{\infty}(X)$ more intrinsically as follows, with $X$ as above. For an extension field $K$ of $k$, the $K$-rational points of $X_{\infty}$ correspond to morphisms of $k$-algebras
\[  \gamma : \Gamma(X,\mathcal{O}_X) \to K[[t]]. \]
For $\lambda\in k^{\times}$ we have an automorphism $\varphi_{\lambda}$ of $K[[t]]$ determined by $t\mapsto \lambda t$. By composing $\varphi_{\lambda}$ with $\gamma$ we obtain a natural action of $k^{\times}$ on $X_{\infty}$ and hence on its coordinate ring $J_{\infty}(X)$. An element $f\in J_{\infty}(X)$ is then called homogeneous of weight $i$ if $\lambda\cdot f = \lambda^i f$, for all $\lambda \in k^{\times}$.  
\end{Rem}

We consider the truncation operator
\[  \tau_{\leq r} : k[[t]] \to k[t] : \sum_{i\geq 0} a_it^i \mapsto \sum_{i=0}^r a_it^i.\] 
Then we have the following simple observation.

\begin{Prop}\label{prop:approximate_by_jet_spaces}
$\tau_{\leq m} \HP_{J_{m}^{\mathfrak{p}}(X)}(t) = \tau_{\leq m} \HP_{J_{\infty}^{\mathfrak{p}}(X)}(t).$
\end{Prop}

Now let $X$ and $Y$ be closed subschemes of $\mathbb{A}^n_k$ respectively $\mathbb{A}^m_k$. Recall that one calls $\mathfrak{p}\in X$ and $\mathfrak{q}\in Y$ \textit{analytically isomorphic} if $\widehat{\mathcal{O}}_{X,\mathfrak{p}}$ and $\widehat{\mathcal{O}}_{Y,\mathfrak{q}}$ are isomorphic $k$-algebras.

\begin{Prop}\label{prop:series_is_analytic_invariant}
If $\mathfrak{p}\in X$ and $\mathfrak{q}\in Y$ are analytically isomorphic then 
\[ \HP_{J_{\infty}^{\mathfrak{p}}(X)}(t) =  \HP_{J_{\infty}^{\mathfrak{q}}(Y)}(t).\]
\end{Prop}

\begin{proof}
The fiber of $\pi_X : X_{\infty} \to X$ over $\pe$ is a scheme over $\kappa(\mathfrak{p})$ whose $K$-rational points, for an extension field $K$ of $\kappa(\mathfrak{p})$, correspond to the set of morphisms of $k$-algebras
\[ \gamma: \Gamma(X,\mathcal{O}_X) \to K[[t]]          \]
such that $\gamma^{-1}\bigl((t)\bigr) = \mathfrak{p}$. Since $K[[t]]$ is complete, $\gamma$ factors uniquely through $\widehat{\mathcal{O}}_{X,\mathfrak{p}}$. We conclude that the fiber of $\pi_X : X_{\infty} \to X$ over $\pe$ is determined by $\widehat{\mathcal{O}}_{X,\mathfrak{p}}$. 

To see that the graded structure of the fibers of $\pi_X$ and $\pi_Y$ above $\mathfrak{p}$ respectively $\mathfrak{q}$ agree, we can use the intrinsic description of the graded structure from the previous remark.
\end{proof}

Next we compute the arc Hilbert-Poincar\'{e} series at a smooth point. We use the notation
\[  \mathbb{H} := \prod_{i\geq 1} \frac{1}{1-t^i}.    \]

\begin{Prop}\label{prop:smooth_case}
Let $X$ be an irreducible closed subscheme of $\mathbb{A}^n_k$ of dimension $d$ and let $\mathfrak{p}\in X$ be a smooth point. Then
\[  \HP_{J_{\infty}^{\mathfrak{p}}(X)}(t) = \mathbb{H}^{d}. \]
\end{Prop}

\begin{proof}
By definition $\mathcal{O}_{X,\mathfrak{p}}$ is a regular local ring, and hence an integral domain. Therefore $X$ is reduced, and thus an integral scheme.
Denote by $e$ the transcendence degree of $\kappa(\mathfrak{p})$ over $k$. From dimension theory (e.g.\ Thm.\ A and Cor.\ 13.4 on p.290 in \cite{eisenbud_commalg}) it follows that the dimension of $\mathcal{O}_{X,\mathfrak{p}}$ equals $d-e$. The complete local ring $\widehat{O}_{X,\mathfrak{p}}$ is regular as well and hence isomorphic to $\kappa(\mathfrak{p})[[x_1,\ldots , x_{d-e}]]$ (Prop.\ 10.16 in \cite{eisenbud_commalg}). By the theorem of the primitive element, $\kappa(\mathfrak{p})$ is isomorphic to $k(y_1,\ldots,y_e)[x]/(f)$, where we may assume $f\in k[y_1,\ldots,y_e,x]$. Hence $\kappa(\mathfrak{p})$ is isomorphic to the residue field of the point $(f,z_1,\ldots , z_{d-e-1})$ in the affine space 
\[ \Spec k[y_1,\ldots,y_e,x,z_1,\ldots, z_{d-e-1}].     \]
From Proposition~\ref{prop:series_is_analytic_invariant} it follows that it suffices to compute the arc Hilbert-Poincar\'{e} series at a point $\mathfrak{q}$ of $Y:=\mathbb{A}^{d}_{k}$. This is an easy task, since $J_{m}(Y)$ is the polynomial ring
\[ k[x_j^{(i)}; 1\leq j \leq d, 0\leq i\leq m],\] 
and $J_{m}^{\mathfrak{q}}(Y)$ equals
\[ \kappa(\mathfrak{q})[x_j^{(i)}; 1\leq j \leq d, 1\leq i\leq m].\] 
The variables form a regular sequence in this ring, and hence it follows easily from Lemma~\ref{lem:hilbert_trick_1} that
\[ \HP_{J_{m}^{\mathfrak{q}}(Y)}(t) = \left(\prod_{i=1}^m \frac{1}{1-t^i} \right)^{d}.\]
Now we use Proposition~\ref{prop:approximate_by_jet_spaces} to finish the proof.
\end{proof}

In Section~\ref{sec:section_4} we discuss the connection of this result with partitions. We leave the proof of the following proposition to the reader.

\begin{Prop}\label{prop:multiplicativity_for_products}
Let $X$ and $Y$ be closed subschemes of $\mathbb{A}^n_k$ respectively $\mathbb{A}^m_k$. Let $\mathfrak{p}\in X$ and $\mathfrak{q}\in Y$. Then
\[ \HP_{J_{\infty}^{(\mathfrak{p},\mathfrak{q})}(X\times Y)}(t) = \HP_{J_{\infty}^{\mathfrak{p}}(X)}(t)\cdot \HP_{J_{\infty}^{\mathfrak{q}}(Y)}(t).  \]
\end{Prop}

The multiplicity of a singular point on a hypersurface can be easily read from the arc Hilbert-Poincar\'{e} series, as the reader may convince himself of:

\begin{Prop}\label{prop:multiplicity_from_HPseries}
Let $X$ be a hypersurface in $\mathbb{A}^n_k$ defined by a polynomial $F\in k[x_1,\ldots,x_n]$ with $F(0)=0$. Then $X$ has multiplicity $r$ at the origin if and only if $r$ is the maximal number such that
\[  \tau_{\leq r-1} \HP_{J_{\infty}^{0}(X)}(t) = \tau_{\leq r-1} \mathbb{H}^n. \]
Moreover, $\tau_{\leq r} \HP_{J_{\infty}^{0}(X)}(t) = \tau_{\leq r} \mathbb{H}^n - t^r$.
\end{Prop}

Next we derive a formula for the arc Hilbert-Poincar\'{e} series of the focussed arc algebra at a canonical hypersurface singularity of maximal multiplicity. First we recall the definition of a canonical singularity. Let $X$ be a normal variety. Assume that $X$ is $\mathbb{Q}$-Gorenstein (i.e.\ $rK_X$ is Cartier for some $r\geq 1$). Let $f: Y\to X$ be a log resolution. This means that $f$ is a proper birational morphism from a smooth variety $Y$ such that the exceptional locus is a simple normal crossings divisor with irreducible components $E_i, i\in I$. We have a linear equivalence
\[  K_Y = f^*K_X + \sum_i a_i E_i \] 
for uniquely determined $a_i\in \mathbb{Q}$ (these are called discrepancy coefficients). Then $X$ has \textit{canonical singularities} if $a_i \geq 0$ for all $i$. We say that $X$ has a canonical singularity at a point $\mathfrak{p}\in X$ if there exists a neighbourhood $U$ of $\mathfrak{p}$ in $X$ with canonical singularities. Note that if $X$ is a hypersurface in $\mathbb{A}^n_k$, then $X$ is Gorenstein (i.e.\ $K_X$ is Cartier) and all $a_i$ are then integers. In that case, if $X$ has a canonical singularity at a closed point $\mathfrak{p}$, the multiplicity of $\mathfrak{p}$ is at most the dimension of $X$. This follows by computing the discrepancy coefficient of the exceptional divisor of the blowing-up in $\mathfrak{p}$. 

\begin{Prop}\label{prop:canonical_sing}
Let $X$ be a normal hypersurface in $\mathbb{A}^n_k$ with a canonical singularity of multiplicity $n-1$ at the origin. Then 
\[ \HP_{J_{\infty}^{0}(X)}(t) = \left(\prod_{i=1}^{n-2} \frac{1}{1-t^i}\right)^{n} \left(\prod_{i\geq n-1} \frac{1}{1-t^i}\right)^{n-1} .\]
\end{Prop}

\begin{proof}
Let $X$ be defined by the polynomial $F$. We use the notations $F_i$ and $f_i$ as before. Then $f_i = 0$ for $0\leq i \leq n-2$. To deduce the result, it suffices to show that for every $m\geq n-1$ the polynomials $f_{n-1}, f_n, \ldots , f_m$ form a regular sequence in the polynomial ring $k[x_j^{(i)};1\leq j \leq n, 1\leq i\leq m]$, in view of Lemma~\ref{lem:hilbert_trick_1} and Proposition~\ref{prop:approximate_by_jet_spaces}. 

Since the question is local, we may assume that all singularities of $X$ are canonical. We will use a theorem by Ein and Musta\c{t}\u{a} that characterizes canonical singularities by the fact that their jet spaces are irreducible (see Thm.\ 1.3 in \cite{ein_mustata}). It is well known that the natural maps $\pi_m : X_m \to X$ are locally trivial fibrations above the smooth part of $X$, with fiber isomorphic to $\mathbb{A}_k^{(n-1)m}$. Hence the dimension of $X_m$ is precisely $(n-1)(m+1)$. Since $X_m$ is irreducible, it follows that the fiber $\pi_m^{-1}(0)$ has dimension at most $(n-1)(m+1)-1$. From dimension theory (e.g.\ Cor.\ 13.4 in \cite{eisenbud_commalg}) we deduce that the codimension of the ideal $I_m^{0}:=(f_{n-1}, f_n, \ldots , f_m)$ in $A_m:=k[x_j^{(i)};1\leq j \leq n, 1\leq i\leq m]$ is at least $nm - \bigl((n-1)(m+1) - 1\bigr) = m-n+2$. From the principal ideal theorem (Thm.\ 10.2 of \cite{eisenbud_commalg}) we get then that the codimension of $I_m^{0}$ is precisely $m-n+2$. Since a polynomial ring over a field is Cohen-Macaulay, we may apply the unmixedness theorem (Cor.\ 18.14 of \cite{eisenbud_commalg}) to deduce that every associated prime of $I_m^{0}$ is minimal. But the codimension of $I_{m+1}^{0}$ in $A_{m+1}$ is at least $m-n+3$, so this means that $f_{m+1}$ is not contained in any minimal prime ideal containing $I_m^{0}$, considered as ideal of $A_{m+1}$. Hence $f_{m+1}$ does not belong to an associated prime ideal of $I_m^{0}$, and thus it is a nonzerodivisor modulo $I_m^{0}$ (see Thm.\ 3.1(b) of \cite{eisenbud_commalg}).
\end{proof}

It follows that the arc Hilbert-Poincar\'{e} series is in this case completely determined by the multiplicity. As a corollary, we get the following nice example. This result was obtained by explicit computation in \cite{mourtada_thesis}.

\begin{Cor}\label{cor:rational_double_point}
If $X$ is a surface with a rational double point at $\mathfrak{p}$ then 
\[ \HP_{J_{\infty}^{\mathfrak{p}}(X)}(t) =  \left(\frac{1}{1-t}\right)^{3} \left(\prod_{i\geq 2} \frac{1}{1-t^i}\right)^2.\]
\end{Cor}

A similar result is true for normal crossings singularities. A scheme $X$ of finite type over $k$ of dimension $d$ is said to have \textit{normal crossings} at a point $\mathfrak{p}$ if $\mathfrak{p}\in X$ is analytically isomorphic to a point $\mathfrak{q}\in Y$, where $Y$ is the hypersurface in $\mathbb{A}_k^{d+1}$ defined by $y_1\cdots y_{d+1} = 0$. For points on $Y$ the situation is as follows.

\begin{Prop}\label{prop:normal_crossings}
Let $Y$ be as above, and assume that $\mathfrak{q}$ lies precisely on the irreducible components given by $y_1=0,\ldots , y_e=0$. Then 
\[ \HP_{J_{\infty}^{\mathfrak{q}}(Y)}(t) =  \left(\prod_{i=1}^{e-1} \frac{1}{1-t^i}\right)^{d+1} \left(\prod_{i\geq e} \frac{1}{1-t^i}\right)^d .\]
\end{Prop}

\begin{proof}
Locally at $\mathfrak{q}$, the variety $Y$ looks like a product of the hypersurface $Z$ given by $z_1\cdots z_e =0$ in $\mathbb{A}^e_k$ and the affine space $\mathbb{A}^{d+1-e}_k$. By Propositions~\ref{prop:multiplicativity_for_products} and \ref{prop:smooth_case} it suffices now to show that 
\begin{equation}\label{formula_normal_crossings}  \HP_{J_{\infty}^{0}(Z)}(t)  =  \left(\prod_{i=1}^{e-1} \frac{1}{1-t^i}\right)^{e} \left(\prod_{i\geq e} \frac{1}{1-t^i}\right)^{e-1}.  \end{equation}
According to Theorem 2.2 of \cite{goward_smith}, the $m$th jet scheme $Z_m$ is equidimensional of dimension $(e-1)(m+1)$, and for $m+1\geq e$ there are actually irreducible components of that dimension in the fiber of $Z_m$ above the origin in $Z$. We can use a reasoning as in the proof of Proposition~\ref{prop:canonical_sing} to conclude that the $m-e+1$ equations $f_e,f_{e+1} , \ldots , f_m$ form a regular sequence in $k[z_j^{(i)}; 1\leq j \leq e, 1\leq i\leq m]$ and then we use Lemma~\ref{lem:hilbert_trick_1} once more to deduce formula~(\ref{formula_normal_crossings}).
\end{proof}